
\chapter{סיכום ועבודה עתידית}

\section{תרומות המחקר}

מחקר זה הציג ארכיטקטורת \textenglish{BiLSTM Encoder-Decoder} לחילוץ גלי סינוס מאותות מעורבבים רועשים. הושגו שיפורים עקביים של \textenglish{2--4 dB} ב-\textenglish{SI-SNR} על פני כלל קווי הבסיס. הממצאים מראים כי:

\begin{itemize}
\item עיבוד דו-כיווני (\textenglish{bidirectional}) מהותי לחילוץ מדויק של פרמטרי תדר.
\item פונקציית אובדן \textenglish{SI-SNR} עדיפה על \textenglish{MSE} בתנאי רעש חזק.
\item \textenglish{Curriculum learning} משפר התכנסות בכ-\textenglish{15\%} ממספר ה-\textenglish{epochs}.
\end{itemize}

\section{מגבלות}

המודל הנוכחי מניח מספר מרכיבים ידוע מראש (\textenglish{K} קבוע). הרחבה לסינוסים עם מספר לא ידוע דורשת גישה של זיהוי (\textenglish{detection}) לפני הערכת הפרמטרים, כמוצע ב-\cite{nachmani2020voice}.

\section{כיוונים עתידיים}

עבודה עתידית תכלול שילוב ה-\textenglish{BiLSTM} עם שכבות קונבולוציה זמניות לחילוץ תכונות מקומיות \cite{luo2019convtasnet}, וכן התאמה לעיבוד בזמן אמת (\textenglish{real-time}) על מחשבי קצה (\textenglish{edge devices}) בעלי זיכרון מוגבל. כמו כן יש לחקור האם ארכיטקטורת \textenglish{xLSTM} \cite{hochreiter1997lstm} עם שערי מטריצה מציגה יתרון נוסף בהפרדת מרכיבים תדריים סמוכים.
